% from ega0-19b-fr.tex
\oldpage[0\textsubscript{IV-1}]{112}

\begin{enumerate}
  \item[(i)] Il existe un anneau de Cohen $W$ tel que $A$ soit isomorphe à un
  anneau quotient d'un anneau de séries formelles
  $W[[T_1,\ldots,T_n]]$ (et en particulier $A$ est isomorphe à un quotient d'un
  anneau local régulier complet (17.3.8)). Si $A$ contient un corps, il est
  isomorphe à un anneau quotient de $k[[T_1,\ldots,T_n]]$.
  \item[(ii)] Supposons en outre $A$ intègre. Alors il existe un sous-anneau
  $B$ de $A$, tel que~: $1^{\mathrm o}$ $B$ est isomorphe à un anneau de séries
  formelles sur un anneau $C$ qui est un corps ou un anneau de Cohen (ce qui
  entraîne que $B$ est un anneau local régulier et complet (17.3.8))~;
  $2^{\mathrm o}$ $B$ a même corps résiduel que $A$ et l'injection $B\to A$ est
  un homomorphisme local~; $3^{\mathrm o}$ $A$ est une $B$-algèbre finie.
\end{enumerate}
\end{theorem}

\begin{proof}
Soit $\mathfrak m$ l'idéal maximal de $A$. Il existe un anneau de Cohen $W$ dont
le corps résiduel est isomorphe à $k$ (19.8.6, (ii))~; on a donc un
homomorphisme local $W\to A/\mathfrak m$, qui par suite se factorise en
$W\xrightarrow{u}A\to A/\mathfrak m$, où $u$ est un homomorphisme local
(19.8.6, (i)). Pour toute famille finie $(x_i)_{1\leq i\leq n}$ d'éléments de
$\mathfrak m$, il existe alors un homomorphisme local
$v:W[[T_1,\ldots,T_n]]\to A$ prolongeant $u$ et tel que $v(T_i)=x_i$ pour tout
$i$ (Bourbaki, \emph{Alg. comm.}, chap.~III, §~4, n\textsuperscript{o}~5,
prop.~6). Lorsque $A$ contient un corps, il contient un corps premier $P$, dont
$k$ est une extension (nécessairement séparable), et par suite $A$ contient un
corps isomorphe à $k$ (19.6.2)~; on peut alors remplacer $W$ par $k$ dans la
définition précédente de $v$.

\medskip
\noindent(i)\enspace
Prenons d'abord pour les $x_i$ un système de générateurs de $\mathfrak m$. Comme
$W$ a même corps résiduel que $A$, et que les classes des $x_i$ dans l'anneau
gradué $\gr_\bullet(A)$ engendrent $\gr_\bullet(A)$ en tant que $k$-algèbre,
$\gr(v):\gr_\bullet(W[[T_1,\ldots,T_n]])\to\gr_\bullet(A)$ est surjectif~; on
en déduit que $v$ est lui-même surjectif (Bourbaki, \emph{Alg. comm.},
chap.~III, §~2, n\textsuperscript{o}~8, cor.~2 du th.~1). Rappelons que le cas
où $A$ contient un corps a déjà été vu et ne figure ici que pour mémoire
(19.6.3).

\medskip
\noindent(ii)\enspace
Si $A$ contient un corps, il contient un corps $k'$ isomorphe à $k$ comme on
l'a vu~; on considère alors un système de paramètres
$(y_j)_{1\leq j\leq m}$ de $A$ (16.3.6), on prend
$B=k[[T_1,\ldots,T_m]]$ et on considère l'homomorphisme local $w:B\to A$ qui
coïncide dans $k$ avec un isomorphisme $k\to k'$ et qui est tel que
$w(T_j)=y_j$ pour $1\leq j\leq m$. Si $A$ ne contient pas de corps, l'unique
homomorphisme $\mathbf Z_{p\mathbf Z}\to A$ (19.8.3) est nécessairement
injectif (sans quoi, comme $A$ est intègre, son noyau serait l'idéal maximal de
$\mathbf Z_{p\mathbf Z}$ et son image isomorphe à un corps)~; en outre, on a
alors $p>0$ par hypothèse, $\mathbf Z_{p\mathbf Z}$ étant un corps si $p=0$.
L'élément $p\cdot1$ de $A$ (identifié à $p$) est non diviseur de zéro dans $A$,
et est contenu dans $\mathfrak m$, donc (16.3.4 et 16.3.7) il existe une famille
$(z_i)_{1\leq i\leq m-1}$ qui, avec $p$, forme un système de paramètres de $A$.
L'anneau de Cohen $W$ considéré au début de la démonstration est alors un
anneau de valuation discrète de corps résiduel $k$, dans lequel $p$ engendre
l'idéal maximal (19.8.5), et l'homomorphisme unique $u:W\to A$ défini au début
applique $p$ sur lui-même. On prend alors $B=W[[T_1,\ldots,T_{m-1}]]$ et on
considère l'homomorphisme local $w:B\to A$ qui coïncide avec $u$ dans $W$ et
est tel que $w(T_i)=z_i$ pour $1\leq i\leq m-1$. Dans les deux cas, si
$\mathfrak n$ est l'idéal maximal de $B$, il est clair que $\mathfrak nA$ est
un idéal de définition de $A$~; comme en outre $B/\mathfrak n$ et
$A/\mathfrak m$ sont isomorphes, $A$ est un $B$-module quasi-fini
($\mathbf 0_{\mathrm I}$, 7.4.4), donc un $A$-module de type fini puisque $B$
est complet et $A$ séparé pour les topologies $\mathfrak n$-préadiques
($\mathbf 0_{\mathrm I}$, 7.4.1). D'autre part,

